\chapter{Adding and Managing Figures}
\label{chap:figures}
\index{figure}\index{graphics!raster}\index{graphics!vector}\index{caption}\index{subfigure}

A figure should let the reader see evidence, a process, or a relationship that
would be harder to grasp in prose. \LaTeX{} does not create most scientific plots;
it places a prepared graphic, supplies a caption and number, and integrates the
figure with the document.

\begin{learningoutcomes}
  \item Distinguish raster and vector graphics and choose suitable formats.
  \item Include an image from a relative path and scale it safely.
  \item Build a numbered figure with caption and label.
  \item Explain float placement without treating movement as an error.
  \item Create a multi-panel figure with \pkg{subcaption}.
  \item Check resolution, colour, grayscale, and caption accessibility.
\end{learningoutcomes}

\section{Prepare the graphic before placing it}

A \term{raster image} stores a grid of pixels. Photographs, microscopy images,
and scans are naturally raster; PNG and JPEG are common formats. A
\term{vector graphic} stores geometric paths and text. Diagrams, line drawings,
and many statistical plots remain sharp when scaled; PDF is the most convenient
vector format for pdf\LaTeX{}.

Use JPEG for continuous-tone photographs when appropriate, PNG for lossless
raster graphics, and PDF for vector artwork. Do not convert a low-resolution
JPEG to PDF and call it vector; the pixels remain.

\begin{goodpractice}
Export plots at their intended publication dimensions with readable labels.
Keep editable plot code or source beside the export. Do not enlarge a small
raster image in \LaTeX{} and expect new detail to appear.
\end{goodpractice}

\section{The smallest included graphic}

Load \pkg{graphicx}, place the file inside the project, and use a relative
path:

\begin{latexcode}{Complete project example: include one PDF graphic}
\documentclass{article}
\usepackage{graphicx}
\begin{document}
\includegraphics[width=0.75\linewidth]{method_comparison.pdf}
\end{document}
\end{latexcode}

This complete source assumes the original file
\file{method\_comparison.pdf} is beside \file{main.tex}. The optional argument
sets the displayed width to 75\% of the current line width; the height changes
proportionally. Do not give both width and height unless distortion is intended.

\section{Caption, label, and float}

Wrap the graphic in a \code{figure} environment:

\begin{figure}[htbp]
  \centering
  \includegraphics[width=0.78\linewidth]
    {figures/original/method_comparison.pdf}
  \caption{Relationship between two methods for simulated paired values. The
  dashed line marks equality; points above it have larger Method B values.}
  \label{fig:method-comparison}
\end{figure}

The source uses \cmd{centering} inside the float, includes the graphic, writes
an informative caption, and labels the number produced by the caption. The
figure is original and the graphic itself states that the data are simulated.

\begin{scienceinpractice}
The dashed line in \cref{fig:method-comparison} provides a reference, but a
cloud near that line is not a complete agreement analysis. A figure must not
imply a stronger scientific conclusion than the method supports.
\end{scienceinpractice}

\section{Why a figure moves}

A float is allowed to move so that \LaTeX{} can avoid large gaps, stranded
captions, and figures cut by page boundaries. The placement letters
\code{h}, \code{t}, \code{b}, and \code{p} mean here, top, bottom, and float
page. \LaTeX{} considers them with size and page constraints.

Place a figure near its first discussion in source and refer to it by label.
If it moves one page, the reference remains correct. Repeatedly adding stronger
placement commands can make the surrounding pages worse.

\begin{commonerror}
\textbf{Symptom:} a figure appears later than its source position.

\textbf{Cause:} insufficient legal space remained on the page, or earlier
floats were waiting.

\textbf{Response:} confirm that the graphic is sensibly sized and placed near
its first mention. Inspect nearby floats. Treat movement as layout behaviour,
not automatically as failure.
\end{commonerror}

\section{Multi-panel figures}

Load \pkg{subcaption}. Each \code{subfigure} gets its own caption and label,
while the outer figure gets the overall caption and number.

\begin{figure}[htbp]
  \centering
  \begin{subfigure}[t]{0.48\textwidth}
    \centering
    \includegraphics[width=\linewidth]
      {figures/original/method_comparison.pdf}
    \caption{Paired values and equality line.}
    \label{fig:multi-values}
  \end{subfigure}\hfill
  \begin{subfigure}[t]{0.48\textwidth}
    \centering
    \includegraphics[width=\linewidth]
      {figures/original/difference_plot.pdf}
    \caption{Difference against pair mean.}
    \label{fig:multi-differences}
  \end{subfigure}
  \caption{Two views of the same simulated teaching data. Panel (a) shows the
  original paired scale; panel (b) makes signed differences visible.}
  \label{fig:multi-panel-comparison}
\end{figure}

Panel labels should be referenced automatically. The overall caption explains
why the panels belong together and defines visual encodings that are not
obvious. Make axis type, units, and sample size visible in the figure or
caption as the target style requires.

\section{Landscape and rotated material}

Large maps or wide tables may need a landscape page. The \pkg{pdflscape}
package provides a \code{landscape} environment and rotates the PDF view for a
reader:

\begin{latexcode}{Optional landscape pattern}
\usepackage{pdflscape}

\begin{landscape}
  % Place one genuinely wide figure or table here.
\end{landscape}
\end{latexcode}

This is a preamble fragment followed by a body fragment. Do not rotate a figure
merely because its labels were exported too small. First redesign the graphic.
Publisher systems may request landscape figures as separate files.

\section{Quality and accessibility checks}

Before submission, inspect the graphic at 100\% PDF view and on a printed test
if print matters. Check:

\begin{itemize}
  \item labels, symbols, and lines at final size;
  \item units and definitions of abbreviations;
  \item raster resolution against current venue instructions;
  \item embedded or outlined fonts where required;
  \item colour combinations in grayscale and for colour-vision differences;
  \item distinction by shape, line style, or label rather than colour alone;
  \item a caption that communicates the figure's purpose and evidence.
\end{itemize}

An informative caption assists all readers, but it is not a complete substitute
for tagged-PDF alternative text in an accessible publishing workflow. Ask the
publisher which accessibility metadata and source they accept.

\begin{submissioncheck}
Open every final exported figure rather than trusting its filename. Confirm
that no crop removed labels, no panel order changed, and the manuscript cites
the correct file. Keep permission and licence records for every non-original
asset.
\end{submissioncheck}

\section{Code lab: a figure that still compiles during drafting}

\begin{codelab}
Place a file named \path{measurement-methods.pdf} beside this source and
compile. Rename the image temporarily and compile again. The framed fallback
keeps the draft buildable while making the missing asset visible.
\end{codelab}

\begin{latexcode}{A complete figure document with a visible fallback}
\documentclass{article}
\usepackage{graphicx}

\begin{document}
\begin{figure}[htbp]
  \centering
  \IfFileExists{measurement-methods.pdf}{%
    \includegraphics[width=.82\linewidth]
      {measurement-methods.pdf}%
  }{%
    \fbox{\rule{0pt}{45mm}\rule{.78\linewidth}{0pt}}%
  }
  \caption{Agreement between two simulated measurement methods.}
  \label{fig:agreement}
\end{figure}

Figure~\ref{fig:agreement} is referenced from the prose.
\end{document}
\end{latexcode}

\chapterproject

\begin{miniproject}
Add the supplied original \file{method\_comparison.pdf} to the running article
with a caption that says the values are simulated. Add the two-panel figure as
a later diagnostic. Label the outer figure and both panels. Mention each figure
from Results without typing a number. Record the editable TikZ sources and
their original status in the asset log.
\end{miniproject}

\chaptersummary

Raster graphics store pixels; vector graphics store scalable geometry.
\pkg{graphicx} includes files from relative paths and scales them
proportionally. A \code{figure} float integrates placement, caption, number,
and label. Movement is normal page design. \pkg{subcaption} organises panels;
\pkg{pdflscape} can support genuinely wide material. Publication checks cover
resolution, labels, units, grayscale, accessibility, provenance, and captions.

\chaptercheck
\begin{chapterchecklist}
  \item The chosen format suits photographic or line-based content.
  \item The file path is relative and portable.
  \item Caption and label follow the included graphic.
  \item Text and lines remain legible at final size.
  \item Meaning does not depend on colour alone.
  \item Original or licensed provenance is recorded.
\end{chapterchecklist}

\chapterexercisestitle
\begin{chapterexercises}
  \item Classify a photograph, logo, line plot, and scanned map as primarily
        raster or vector.
  \item Include a PDF graphic at half the line width.
  \item Add a caption and correctly positioned label.
  \item Explain why the \optional{h} request may not place a float exactly here.
  \item Build two subfigures with an overall caption.
  \item Identify two ways to preserve meaning in grayscale.
  \item Diagnose a ``file not found'' error caused by a machine-specific path.
  \item Challenge: audit one figure for evidence, resolution, accessibility,
        caption completeness, and provenance.
\end{chapterexercises}
